% TODO mettre des flèches entre les rectangles rouges pour bien indiquer
% le sens de l’empilement
\begin{tikzpicture}[
        every node/.style={rectangle, font=\footnotesize, inner sep=3pt},
        every state/.style={rectangle, font=\footnotesize, inner sep=1pt},
        every loop/.style={min distance=1mm, looseness=2},
        auto,
        node distance=7mm and 7mm
    ]

    \node (q0) [state, initial, initial text={}] {
        $\begin{aligned}
            \underline{P} & \to \pointeur \text{V égal E} \\
            V             & \to \pointeur \text{id}
        \end{aligned}$
    };

    \uncover<6-34>{
        \node[draw, rectangle, red, thick] at (q0) {
            \phantom{$\begin{aligned}
                \underline{P} & \to \pointeur \text{V égal E} \\
                V             & \to \pointeur \text{id}
            \end{aligned}$}
        };
    }

    \node (q1) [state, above=of q0] {
        $ V \to \text{id} \pointeur $
    };

    \uncover<7>{
        \node[draw, rectangle, red, thick, inner ysep=10pt] at (q1) {
            \phantom{$ V \to \text{id} \pointeur $}
        };
    }

    \node (q2) [state, below=of q0] {
        $ \underline{P} \to \text{V \pointeur égal E} $
    };

    \uncover<9-33>{
        \node[draw, rectangle, red, thick, inner ysep=10pt] at (q2) {
            \phantom{ $ \underline{P} \to \text{V \pointeur égal E} $ }
        };
    }

    \node (q3) [state, below=of q2] {
        $\begin{aligned}
            \underline{P} & \to \text{V égal \pointeur E} \\
            E             & \to \pointeur \text{ent} \\
            E             & \to \pointeur \text{parg E O E pard}
        \end{aligned}$
    };

    \uncover<10-33>{
        \node[draw, rectangle, red, thick] at (q3) {
            \phantom{ $\begin{aligned}
                \underline{P} & \to \text{V égal \pointeur E} \\
                E             & \to \pointeur \text{ent} \\
                E             & \to \pointeur \text{parg E O E pard}
            \end{aligned}$ }
        };
    }

    \node (q3bis) [state, below=of q3] {
        $ \underline{P} \to \text{V égal E \pointeur} $
    };

    \uncover<33>{
        \node[draw, rectangle, red, thick, inner ysep=10pt] at (q3bis) {
            \phantom{$ \underline{P} \to \text{V égal E \pointeur} $}
        };
    }

    \node (q4) [state, right=of q3] {
        $ E \to \text{ent} \pointeur $
    };

    \uncover<12, 19, 25>{
        \node[draw, rectangle, red, thick, inner ysep=10pt] at (q4) {
            \phantom{ $ E \to \text{ent} \pointeur $ }
        };
    }

    \node (q5) [state, above right=of q3] {
        $\begin{aligned}
            E & \to \text{parg \pointeur E O E pard} \\
            E & \to \text{\pointeur ent} \\
            E & \to \text{\pointeur parg E O E pard}
        \end{aligned}$
    };

    \uncover<11-31>{
        \node[draw, rectangle, red, thick] at (q5) {
            \phantom{
            $\begin{aligned}
                E & \to \text{parg \pointeur E O E pard} \\
                E & \to \text{\pointeur ent} \\
                E & \to \text{\pointeur parg E O E pard}
            \end{aligned}$
            }
        };
    }

    \uncover<18-28>{
        \node[draw, rectangle, red, thick, inner sep=6pt] at (q5) {
            \phantom{
            $\begin{aligned}
                E & \to \text{parg \pointeur E O E pard} \\
                E & \to \text{\pointeur ent} \\
                E & \to \text{\pointeur parg E O E pard}
            \end{aligned}$
            }
        };
    }

    \node (q6) [state, above=of q5] {
        $\begin{aligned}
            E & \to \text{parg E \pointeur O E pard} \\
            O & \to \pointeur \text{plus} \\
            O & \to \pointeur \text{fois}
        \end{aligned}$
    };

    \uncover<14-31>{
        \node[draw, rectangle, red, thick] at (q6) {
            \phantom{
            $\begin{aligned}
                E & \to \text{parg E \pointeur O E pard} \\
                O & \to \pointeur \text{plus} \\
                O & \to \pointeur \text{fois}
            \end{aligned}$
            }
        };
    }

    \uncover<21-28>{
        \node[draw, rectangle, red, thick, inner sep=6pt] at (q6) {
            \phantom{
            $\begin{aligned}
                E & \to \text{parg E \pointeur O E pard} \\
                O & \to \pointeur \text{plus} \\
                O & \to \pointeur \text{fois}
            \end{aligned}$
            }
        };
    }


    \node (q7) [state, above left=of q6, xshift=15mm] {
        $ O \to \text{plus} \pointeur $
    };
    \node (q8) [state, above right=of q6, xshift=-15mm] {
        $ O \to \text{fois} \pointeur $
    };

    \uncover<22>{
        \node[draw, rectangle, red, thick, inner ysep=10pt] at (q7) {
            \phantom{$ O \to \text{plus} \pointeur $}
        };
    }
    \uncover<15>{
        \node[draw, rectangle, red, thick, inner ysep=10pt] at (q8) {
            \phantom{$ O \to \text{fois} \pointeur $}
        };
    }

    \node (q9) [state, right=of q6] {
        $\begin{aligned}
            E & \to \text{parg E O \pointeur E pard} \\
            E & \to \pointeur \text{ent} \\
            E & \to \pointeur \text{parg E O E pard}
        \end{aligned}$
    };

    \uncover<17-31>{
        \node[draw, rectangle, red, thick] at (q9) {
            \phantom{
            $\begin{aligned}
                E & \to \text{parg E O \pointeur E pard} \\
                E & \to \pointeur \text{ent} \\
                E & \to \pointeur \text{parg E O E pard}
            \end{aligned}$
            }
        };
    }
    \uncover<24-28>{
        \node[draw, rectangle, red, thick, inner sep=6pt] at (q9) {
            \phantom{
            $\begin{aligned}
                E & \to \text{parg E O \pointeur E pard} \\
                E & \to \pointeur \text{ent} \\
                E & \to \pointeur \text{parg E O E pard}
            \end{aligned}$
            }
        };
    }

    \node (q10) [state, below right=of q9, xshift=-25mm] {
        $ E \to \text{parg E O E \pointeur pard} $
    };
    \node (q11) [state, below=of q10] {
        $ E \to \text{parg E O E pard \pointeur} $
    };

    \uncover<27-28, 30-31>{
        \node[draw, rectangle, red, thick, inner ysep=10pt] at (q10) {
            \phantom{$ E \to \text{parg E O E \pointeur pard} $}
        };
    }
    \uncover<28, 31>{
        \node[draw, rectangle, red, thick, inner ysep=10pt] at (q11) {
            \phantom{$ E \to \text{parg E O E pard \pointeur} $}
        };
    }


    \path[->, every node/.style={font=\footnotesize}]
        (q0)  edge              node [swap]                    {id}    (q1)
              edge              node                           {V}     (q2)
        (q2)  edge              node                           {égal}  (q3)
        (q3)  edge              node                           {ent}   (q4)
              edge              node                           {parg}  (q5)
              edge              node                           {E}     (q3bis)
        (q5)  edge [loop right] node [yshift=4mm, xshift=-1mm] {parg}  ()
              edge              node                           {ent}   (q4)
              edge              node                           {E}     (q6)
        (q6)  edge              node [yshift=2.5mm]            {plus}  (q7)
              edge              node [yshift=-2.5mm]           {fois}  (q8)
              edge              node                           {O}     (q9)
        (q9)  edge              node [auto=right, yshift=-1mm] {parg}  (q5)
              edge [bend left]  node                           {ent}   (q4)
              edge              node [yshift=-2.1mm]           {E}     (q10)
        (q10) edge              node                           {pard}  (q11);

\end{tikzpicture}
